% =============================================================================
% 6. Business Analysis and Decision Framework -- target ~4.5 pp
% THE BA-CORE OF THE THESIS. This is what distinguishes a Business Analytics
% thesis from a generic ML report.
% Three sub-themes:
%   1. Why preference alignment is becoming a strategic capability
%   2. A decision framework for choosing a training method
%   3. The alignment economics that drive the framework
% =============================================================================

\section{Business Analysis and Decision Framework}
\label{sec:business}

\subsection{Market Tendencies}
\label{sec:business:tendencies}

When the first large language models started getting widespread attention, there was a common belief that model performance came mainly from how much data you trained on \citep{brown2020gpt3}. It quickly became apparent that this was not the whole story, and that aligning these models with what users actually wanted was what made them useful in practice. The first widely adopted alignment framework was Reinforcement Learning from Human Feedback, or RLHF, originally proposed by \citet{christiano2017rlhf} for game and robotics environments. The recipe is simple in principle: humans compare pairs of model outputs, a reward model is trained to predict those preferences, and the language model is then optimised against that learned reward. \citet{ouyang2022instructgpt} scaled the recipe to GPT-3 and produced InstructGPT, the direct ancestor of ChatGPT and every aligned LLM that followed. The leap in usefulness was enormous, but the architecture was still one-size-fits-all. A single reward model represented the preferences of a small pool of human annotators, and the resulting policy was shipped to every user as if their tastes matched that pool's.

After this generation of systems reached mass adoption, it became clear that the next step was aligning the model with each individual user. The way the industry has approached this is through context. Instead of retraining the model for every user, you feed it information about who is asking, what they care about, or what task they are trying to do, and let the model condition on that text at inference time. The underlying technique was popularised by Retrieval-Augmented Generation \citep{lewis2020rag} and scaled by follow-up work such as RETRO \citep{borgeaud2022retro}. Productised versions show up everywhere today, from Claude Projects and ChatGPT Memory to Gemini's million-token context window \citep{reid2024gemini}. The appeal is clear: the model keeps all the general world knowledge it learned during pretraining and gets steered toward the task at hand without any weight updates. The important caveat, documented by \citet{liu2023lostmiddle}, is that more context is not always better. Models do not attend uniformly to long prompts, and information buried in the middle of the window is often ignored. This matters for personalisation specifically, because user preferences are exactly the kind of information that piles up over time and risks getting drowned out.

Recommendation systems have always been central to how companies understand their users and connect them with the products they are most likely to want, and for most of their history they have lived behind an interface of ranked lists, grids, and carousels. That interface is now shifting. Users are increasingly comfortable asking for what they want in natural language, and the front door to a product is more and more often a conversation rather than a feed. The academic side has followed quickly. \citet{bao2023tallrec} showed that an LLM can be tuned to act as the recommender itself, \citet{friedman2023recllm} demonstrated this in an explicitly conversational setting, and the survey by \citet{fan2024recllmsurvey} maps how the broader field is reorganising around this transition. The underlying data has not changed, the years of clicks, ratings, and purchases that classical recommenders were built on are still the most valuable signal a company has about its users. What has changed is the surface on which that signal gets used, and the kind of model that needs to consume it.

The findings in this thesis come to inform organisations on how they can tackle these changes and leverage their data, or lack thereof, to enter the LLM world.

\subsection{Business Applications}
\label{sec:business:framework}

These findings become relevant in a business context mainly for companies that use LLMs for customer-facing sales and recommendation, and the natural way to see why is to think about what those companies already have. Almost every consumer business that has been operating for more than a few years sits on a large quantity of discrete user signal, the kind of data classical recommender systems were built to consume. Ratings, watch histories, purchase sequences, thumbs-ups and skips are all the same shape as the preference pairs used throughout this thesis. The work of \citet{bao2023tallrec} and \citet{friedman2023recllm} shows that this signal can be fed into an LLM either by training on it or by placing it in the prompt, and the survey by \citet{fan2024recllmsurvey} confirms this is the direction the field is moving. The question for a company is no longer whether to do it, but how.

That decision is exactly what the results in \cref{sec:results} inform. The crossover between ICL and the trained methods, visible in \cref{fig:primary_curves} and quantified in \cref{tab:primary_acc}, is not just an academic curiosity, it tells a product team where to put a line in their onboarding flow. A music app that hopes a new user will rate ten songs in their first session is operating squarely in the cold-start floor, where ICL leads on both accuracy and cost. A bookseller that nudges every shopper toward fifty ratings before personalising is in the territory where SimPO and IPO have already overtaken the prompt-stuffing baseline. The same business decision, namely how much friction to add to onboarding in exchange for richer personalisation downstream, is also the technical decision of which method to deploy. The two are the same lever seen from two sides.

What makes this practically attractive is that the infrastructure does not change across the regimes. The signal a company needs to collect is the same paired or unpaired preference data either way, and the serving stack only diverges at the last step, with the trained methods loading a small LoRA adapter on top of the same frozen base model that ICL queries directly. \citet{sheng2024slora} showed that thousands of LoRA adapters can be served concurrently on shared infrastructure, which removes the operational objection that per-user training is impractical at scale. The broader point is that the LLM era does not retire the work companies have done on their recommendation stacks over the past decade, it relocates it. The years of clicks and ratings that fed a matrix-factorisation model do not become irrelevant when the front door turns into a chat box, they become the cold-start fuel for whatever personalisation method sits behind that chat box.

\subsection{Alignment Economics}
\label{sec:business:economics}

The decision of whether a company should use in-context learning or shift to a preference-optimisation method is not a one-time architectural choice. It depends on three variables that differ widely between products: how many preferences a user generates during onboarding, how often that same user queries the system afterwards, and how stable their preferences are over time. Mapping the experimental findings from \cref{sec:results} against the cost trajectories documented in the Stanford AI Index 2025 \citep{maslej2025aiindex} lets us turn the abstract trade-off into a decision a product team can actually make.

The empirical anchor is the crossover in \cref{fig:primary_curves}. For $n$ below roughly ten preference pairs, ICL outperforms every trained method on both GoodReads and Netflix, often by a wide margin. Between $n = 10$ and $n = 50$, the trained methods catch up and pass the ICL plateau. At $n = 100$, SimPO leads \texttt{icl\_chat} by 3.6 percentage points on GoodReads and 5.9 points on Netflix, with the gap widening further at $n = 150$. This gives three operational regimes, summarised in \cref{tab:business_regimes}: a cold-start floor where ICL wins on both accuracy and cost, a transition band where the choice is ambiguous, and a steady state where trained methods dominate on accuracy and the question shifts to whether they also win on cost.

\begin{table}[h]
\centering
\caption{Decision regime by training size per user. Recommendations follow directly from the results in \cref{tab:primary_acc} combined with the cost analysis below.}
\label{tab:business_regimes}
\begin{tabular}{lll}
\toprule
\textbf{Regime} & \textbf{$n$ per user} & \textbf{Recommended method} \\
\midrule
Cold-start floor & $n \leq 10$       & ICL \\
Transition band  & $10 < n \leq 50$  & Project engagement before deciding \\
Steady state     & $n > 50$          & SimPO (or IPO with a reference anchor) \\
\bottomrule
\end{tabular}
\end{table}

The Stanford report changes how heavy the cost side of that question is. The cost of running an inference at GPT-3.5 quality dropped from \$20.00 per million tokens in November 2022 to \$0.07 per million tokens by October 2024, a 280-fold reduction in roughly eighteen months \citep{maslej2025aiindex}. Hardware price-performance is improving at 30 percent per year, and a 3.8-billion-parameter Phi-3-mini now matches the MMLU score that a 540-billion-parameter PaLM achieved two years earlier. The practical effect is that the per-query cost of an ICL prompt, which used to be the dominant economic argument against the method, has collapsed. An ICL inference at $n = 100$ carries roughly four thousand extra prompt tokens compared to a trained-model inference. At the Stanford lower bound that adds less than three hundredths of a cent per query. At a more realistic mid-market rate of \$0.20 per million input tokens it is about eight hundredths of a cent. A LoRA training run for one user on a rented A100 costs somewhere between ten and twenty cents, which puts the break-even point at roughly two hundred and fifty queries per user without prompt caching, and around two thousand five hundred queries once prefix caching is in play.\footnote{Anthropic's prompt-caching documentation reports a roughly tenfold cost reduction on cached prefix tokens: \url{https://platform.claude.com/docs/en/build-with-claude/prompt-caching}.}

Two trajectory effects make this picture less stable than \cref{tab:business_regimes} suggests. The first is that ICL economics are improving faster than training economics. \citet{maslej2025aiindex} reports inference cost falling between 9 and 900 times per year depending on the task, while hardware price-performance drops a steadier 30 percent per year. Over the next few cycles, ICL gets cheaper than training does, which pushes the break-even point higher and rewards companies that delay the decision to train. The second is that smaller models are catching larger ones on standard benchmarks, which lifts the floor of base-model capability that ICL exploits. The frozen Llama-3-8B-Instruct used in our experiments is already capable enough to extract signal from book titles at $n = 10$. Two model generations from now, that capability will sit lower on the parameter curve and run faster, which favours ICL further at the entry point.

The recommendation that comes out of this is sequencing rather than choosing. Most consumer products are better off starting with ICL and switching to a trained adapter only for users who cross a clear engagement threshold. The data infrastructure for both is the same, a stream of paired or unpaired preferences accumulated per user. The decision of when to train is then a business question about engagement rather than a technical question about model behaviour, and the empirical crossover in \cref{sec:results} is what tells you where to put the threshold. A company that sees most of its users churn before $n = 10$ should not be training anything. A company that sees its average user generate a hundred preferences a month should be training adapters for everyone past the first session. The Stanford trajectory is what makes this sequenced strategy increasingly attractive, because the cost of getting it wrong in the ICL direction is dropping faster than the cost of getting it wrong in the training direction.
